% GENERATED FILE. Edit canonical lessons, then run npm run generate:books.
% canonical-source-hash: fnv1a64:0526c04eeb91cd0d
% canonical-lessons: MR-C250-itihas, MR-C250-bhugol, MR-C250-ganit, MR-C250-kala, MR-C250-sangh, MR-R250-first-pass-words-from-chapters-246-247, MR-R250-second-pass-words-from-chapters-248-250

\chapter{History, Geography, Mathematics, Art, Team}
\label{ch:mr-history-geography-mathematics-art-team}
\hlchaptermodality{250}

\begin{chapteropening}
\textbf{By the end of this chapter:} \emph{I can say history, geography, mathematics, art and a team.}

Complete the last lesson of chapter 250: I can say history, geography, mathematics, art and a team.
\end{chapteropening}

\section[\texorpdfstring{\mr{इतिहास} (itihās)}{इतिहास (itihās)}]{\mr{इतिहास} (itihās) — history}

\label{lesson:MR-C250-itihas}



\begin{quote}
Before the new one: say the Marathi for a kite, then the Marathi for a spinning top.
\end{quote}

\subsection*{You'll want to know: \mr{इतिहास}}
\textbf{\mr{इतिहास}} — \emph{itihās} — \textquotedblleft{}history\textquotedblright{}.

\begin{grammarlens}[title={What you've built}]
One more word for this chapter.
\end{grammarlens}

\subsection*{Your turn}
\begin{itemize}
\raggedright
  \item \emph{Say it:} \emph{itihās}
  \item \emph{Say it:} \emph{itihās}, once more
  \item \emph{Say it:} say \emph{itihās}
  \item \emph{Recall:} say the Marathi for a kite, then the Marathi for a spinning top, then say \emph{itihās} again
\end{itemize}

\begin{tcolorbox}[breakable,colback=teal!4,colframe=teal!35!black,title={Before you move on}]
What does \mr{इतिहास} mean? (\textquotedblleft{}History\textquotedblright{}.) Say it once more.
\end{tcolorbox}
\section[\texorpdfstring{\mr{भूगोल} (bhūgol)}{भूगोल (bhūgol)}]{\mr{भूगोल} (bhūgol) — geography}

\label{lesson:MR-C250-bhugol}



\begin{quote}
Before the new one: say the Marathi for a spinning top, then the Marathi for history.
\end{quote}

\subsection*{You'll want to know: \mr{भूगोल}}
\textbf{\mr{भूगोल}} — \emph{bhūgol} — \textquotedblleft{}geography\textquotedblright{}.

\begin{grammarlens}[title={What you've built}]
One more word for this chapter.
\end{grammarlens}

\subsection*{Your turn}
\begin{itemize}
\raggedright
  \item \emph{Say it:} \emph{bhūgol}
  \item \emph{Say it:} \emph{bhūgol}, once more
  \item \emph{Say it:} say \emph{bhūgol}
  \item \emph{Recall:} say the Marathi for a spinning top, then the Marathi for history, then say \emph{bhūgol} again
\end{itemize}

\begin{tcolorbox}[breakable,colback=teal!4,colframe=teal!35!black,title={Before you move on}]
What does \mr{भूगोल} mean? (\textquotedblleft{}Geography\textquotedblright{}.) Say it once more.
\end{tcolorbox}
\section[\texorpdfstring{\mr{गणित} (gaṇit)}{गणित (gaṇit)}]{\mr{गणित} (gaṇit) — mathematics}

\label{lesson:MR-C250-ganit}



\begin{quote}
Before the new one: say the Marathi for history, then the Marathi for geography.
\end{quote}

\subsection*{You'll want to know: \mr{गणित}}
\textbf{\mr{गणित}} — \emph{gaṇit} — \textquotedblleft{}mathematics\textquotedblright{}.

\begin{grammarlens}[title={What you've built}]
One more word for this chapter.
\end{grammarlens}

\subsection*{Your turn}
\begin{itemize}
\raggedright
  \item \emph{Say it:} \emph{gaṇit}
  \item \emph{Say it:} \emph{gaṇit}, once more
  \item \emph{Say it:} say \emph{gaṇit}
  \item \emph{Recall:} say the Marathi for history, then the Marathi for geography, then say \emph{gaṇit} again
\end{itemize}

\begin{tcolorbox}[breakable,colback=teal!4,colframe=teal!35!black,title={Before you move on}]
What does \mr{गणित} mean? (\textquotedblleft{}Mathematics\textquotedblright{}.) Say it once more.
\end{tcolorbox}
\section[\texorpdfstring{\mr{कला} (kalā)}{कला (kalā)}]{\mr{कला} (kalā) — art}

\label{lesson:MR-C250-kala}



\begin{quote}
Before the new one: say the Marathi for geography, then the Marathi for mathematics.
\end{quote}

\subsection*{You'll want to know: \mr{कला}}
\textbf{\mr{कला}} — \emph{kalā} — \textquotedblleft{}art\textquotedblright{}.

\begin{grammarlens}[title={What you've built}]
One more word for this chapter.
\end{grammarlens}

\subsection*{Your turn}
\begin{itemize}
\raggedright
  \item \emph{Say it:} \emph{kalā}
  \item \emph{Say it:} \emph{kalā}, once more
  \item \emph{Say it:} say \emph{kalā}
  \item \emph{Recall:} say the Marathi for geography, then the Marathi for mathematics, then say \emph{kalā} again
\end{itemize}

\begin{tcolorbox}[breakable,colback=teal!4,colframe=teal!35!black,title={Before you move on}]
What does \mr{कला} mean? (\textquotedblleft{}Art\textquotedblright{}.) Say it once more.
\end{tcolorbox}
\section[\texorpdfstring{\mr{संघ} (saṅgh)}{संघ (saṅgh)}]{\mr{संघ} (saṅgh) — a team}

\label{lesson:MR-C250-sangh}



\begin{quote}
Before the new one: say the Marathi for mathematics, then the Marathi for art.
\end{quote}

\subsection*{You'll want to know: \mr{संघ}}
\textbf{\mr{संघ}} — \emph{saṅgh} — \textquotedblleft{}a team\textquotedblright{}.

\begin{grammarlens}[title={What you've built}]
One more word for this chapter.
\end{grammarlens}

\subsection*{Your turn}
\begin{itemize}
\raggedright
  \item \emph{Say it:} \emph{saṅgh}
  \item \emph{Say it:} \emph{saṅgh}, once more
  \item \emph{Say it:} say \emph{saṅgh}
  \item \emph{Recall:} say the Marathi for mathematics, then the Marathi for art, then say \emph{saṅgh} again
\end{itemize}

\begin{tcolorbox}[breakable,colback=teal!4,colframe=teal!35!black,title={Before you move on}]
What does \mr{संघ} mean? (\textquotedblleft{}A team\textquotedblright{}.) Say it once more.
\end{tcolorbox}
\section[(dialogue)]{Words from chapters 246-247 — a first pass}

\label{lesson:MR-R250-first-pass-words-from-chapters-246-247}



\begin{quote}
Say the words for art and a team.
\end{quote}

\subsection*{Your turn}
Say these aloud:

\begin{itemize}
\raggedright
  \item a rainbow, hail, springtime, a planet, freedom
  \item a job, a degree, homework, education, a feeling
\end{itemize}

\begin{tcolorbox}[breakable,colback=teal!4,colframe=teal!35!black,title={Before you move on}]
A rainbow? (\textbf{\mr{इंद्रधनुष्य}}, \emph{indradhanuṣya}.) Hail? (\textbf{\mr{गारा}}, \emph{gārā}.) Springtime? (\textbf{\mr{वसंत}}, \emph{vasant}.) A planet? (\textbf{\mr{ग्रह}}, \emph{grah}.) Freedom? (\textbf{\mr{स्वातंत्र्य}}, \emph{svātantrya}.) A job? (\textbf{\mr{नोकरी}}, \emph{nokrī}.) A degree? (\textbf{\mr{पदवी}}, \emph{padvī}.) Homework? (\textbf{\mr{गृहपाठ}}, \emph{gṛhapāṭh}.) Education? (\textbf{\mr{शिक्षण}}, \emph{śikṣaṇ}.) A feeling? (\textbf{\mr{भावना}}, \emph{bhāvanā}.)
\end{tcolorbox}
\section[(dialogue)]{Words from chapters 248-250 — a second pass}

\label{lesson:MR-R250-second-pass-words-from-chapters-248-250}



\begin{quote}
Say the words for a festival and a procession.
\end{quote}

\subsection*{Your turn}
Say these aloud:

\begin{itemize}
\raggedright
  \item a festival, a procession, a village fair, a feast, a gift
  \item a swing, toys, a doll, a kite, a spinning top
  \item history, geography, mathematics, art, a team
\end{itemize}

\begin{tcolorbox}[breakable,colback=teal!4,colframe=teal!35!black,title={Before you move on}]
A festival? (\textbf{\mr{सण}}, \emph{saṇ}.) A procession? (\textbf{\mr{मिरवणूक}}, \emph{miravṇūk}.) A village fair? (\textbf{\mr{जत्रा}}, \emph{jatrā}.) A feast? (\textbf{\mr{मेजवानी}}, \emph{mejvānī}.) A gift? (\textbf{\mr{भेटवस्तू}}, \emph{bheṭvastū}.) A swing? (\textbf{\mr{झोपाळा}}, \emph{jhopāḷā}.) Toys? (\textbf{\mr{खेळणी}}, \emph{kheḷṇī}.) A doll? (\textbf{\mr{बाहुली}}, \emph{bāhulī}.) A kite? (\textbf{\mr{पतंग}}, \emph{pataṅg}.) A spinning top? (\textbf{\mr{भोवरा}}, \emph{bhovrā}.) History? (\textbf{\mr{इतिहास}}, \emph{itihās}.) Geography? (\textbf{\mr{भूगोल}}, \emph{bhūgol}.) Mathematics? (\textbf{\mr{गणित}}, \emph{gaṇit}.) Art? (\textbf{\mr{कला}}, \emph{kalā}.) A team? (\textbf{\mr{संघ}}, \emph{saṅgh}.)
\end{tcolorbox}
